\documentclass[11pt,a4paper]{article}
\usepackage[utf8]{inputenc}
\usepackage{amsmath,amsfonts,amssymb}
\usepackage{graphicx}
\usepackage{geometry}
\geometry{margin=1in}
\usepackage{booktabs}
\usepackage{hyperref}
\usepackage{float}
\usepackage{xcolor}

\definecolor{darkbg}{HTML}{0d121f}
\definecolor{lighttext}{HTML}{e2e8f0}
\definecolor{primary}{HTML}{8b5cf6}
\definecolor{secondary}{HTML}{0dd5c5}

\pagecolor{darkbg}
\color{lighttext}

\hypersetup{
    colorlinks=true,
    linkcolor=secondary,
    urlcolor=secondary
}

\title{\textbf{EE3621: Digital Signal Processing} \\ \Large Lecture 21: Quantization Effects in DSP Systems \\ \large (III B.Tech EEE)}
\author{Lecture Notes (40-Minute Plan)}
\date{}

\begin{document}

\maketitle

\section{Lecture Plan \& Objectives (40 Minutes Breakdown)}
\begin{itemize}
    \item \textbf{00:00 -- 05:00 (5 mins):} Why Quantization Matters (ADC limitations, fixed-point vs floating-point).
    \item \textbf{05:00 -- 12:00 (7 mins):} Uniform Quantization \& Error Model (step size, representation levels, white noise assumption).
    \item \textbf{12:00 -- 18:00 (6 mins):} SQNR Derivation ($6.02b + 1.76$ dB).
    \item \textbf{18:00 -- 25:00 (7 mins):} Coefficient Quantization (FIR/IIR, sensitivity, direct vs cascade forms).
    \item \textbf{25:00 -- 30:00 (5 mins):} Round-off Noise in Fixed-Point IIR (additive noise model, PSD, noise power).
    \item \textbf{30:00 -- 35:00 (5 mins):} Limit Cycles \& Overflow (granular, overflow, saturation vs two's complement).
    \item \textbf{35:00 -- 40:00 (5 mins):} Checkpoints \& Full Derivations.
\end{itemize}

\noindent\rule{\textwidth}{0.5pt}

\section{Why Quantization Matters}
In any practical Digital Signal Processing (DSP) system, we face finite word-length effects due to hardware limitations.
An Analog-to-Digital Converter (ADC) maps continuous amplitudes into a finite set of digital levels, inherently discarding information.
Fixed-point DSPs assign a strict number of bits to the integer and fractional parts. We discard bits after multiplication via truncation (chopping lower bits) or rounding (choosing the nearest representable value).

\section{Uniform Quantization}
Let the input signal range be from $x_{min}$ to $x_{max}$. For a quantizer with $b$ bits, the number of levels is $L = 2^b$. The step size $\Delta$ is:
\begin{equation}
\Delta = \frac{x_{max} - x_{min}}{2^b}
\end{equation}

In the \textbf{Granular Noise Region}, the signal stays within bounds and error is constrained by $\Delta$. In the \textbf{Overload Region}, the signal exceeds bounds and clipping causes severe distortion.

\section{Quantization Error Model}
The quantized signal $x_Q[n]$ is modeled as $x[n]$ plus additive noise $e[n]$:
\begin{equation}
x_Q[n] = x[n] + e[n] \quad \Rightarrow \quad e[n] = x_Q[n] - x[n]
\end{equation}
Under the white noise model (Widrow assumption), $e[n]$ is uniformly distributed over $[-\Delta/2, \Delta/2]$. The variance is:
\begin{align}
\sigma_e^2 &= E[e^2] \\
&= \int_{-\Delta/2}^{\Delta/2} e^2 \frac{1}{\Delta} de \\
&= \frac{1}{\Delta} \left[ \frac{e^3}{3} \right]_{-\Delta/2}^{\Delta/2} \\
&= \frac{\Delta^2}{12}
\end{align}

\section{Signal-to-Quantization-Noise Ratio (SQNR)}
For a full-scale sinusoid $x[n] = A \cos(\omega n)$, peak-to-peak is $2A$. Thus, $\Delta = \frac{2A}{2^b}$.
\begin{align}
P_x &= \frac{A^2}{2} \\
P_e &= \frac{\Delta^2}{12} = \frac{1}{12} \left( \frac{2A}{2^b} \right)^2 = \frac{A^2}{3 \cdot 2^{2b}} \\
\text{SQNR} &= \frac{P_x}{P_e} = \frac{3}{2} \cdot 2^{2b} \\
\text{SQNR (dB)} &= 10 \log_{10} \left( \frac{3}{2} \cdot 2^{2b} \right) = 6.02b + 1.76
\end{align}
Every additional bit increases SQNR by roughly 6 dB.

\section{Coefficient Quantization in FIR/IIR}
Filter coefficients $h_k$ are quantized, perturbing poles and zeros:
\begin{equation}
\Delta H(e^{j\omega}) \approx \sum_k \frac{\partial H(e^{j\omega})}{\partial h_k} \Delta h_k
\end{equation}
High-order direct form IIR filters are highly sensitive to coefficient quantization and can become unstable. Cascade forms (second-order sections) are much less sensitive since root dependency is decoupled.

\section{Round-off Noise in Fixed-Point IIR}
A $b$-bit times $b$-bit multiplication gives $2b$ bits, which is rounded. We model this as inserting white noise $e[n]$ after each multiplier. The output PSD is:
\begin{equation}
S_{ee}^{out}(e^{j\omega}) = \frac{\sigma_e^2}{2\pi} |H_{noise}(e^{j\omega})|^2
\end{equation}
Total output noise power is:
\begin{equation}
\sigma_{out}^2 = \sigma_e^2 \sum_{n=0}^{\infty} h_{noise}^2[n]
\end{equation}

\section{Limit Cycles}
Fixed-point IIR filters can exhibit sustained oscillations for zero input.
\begin{itemize}
    \item \textbf{Granular Limit Cycles:} Rounding/truncation inside the feedback loop.
    \item \textbf{Overflow Limit Cycles:} Occurs due to Two's Complement wraparound. Saturation arithmetic clips the value instead, heavily reducing this risk.
\end{itemize}

\section{Floating-Point vs Fixed-Point}
Fixed-point has constant absolute error. IEEE 754 Floating-point uses a mantissa and exponent ($x = M \cdot 2^E$), offering bounded relative error and a huge dynamic range. It is preferred when scaling is too complex or dynamic range fluctuates drastically.

\section{Key Formulas Summary}
\begin{table}[H]
\centering
\begin{tabular}{ll}
\toprule
\textbf{Concept} & \textbf{Formula} \\
\midrule
Step Size & $\Delta = \frac{x_{max} - x_{min}}{2^b}$ \\
Quantization Noise Variance & $\sigma_e^2 = \frac{\Delta^2}{12}$ \\
SQNR for Full-Scale Sine & $\text{SQNR (dB)} = 6.02b + 1.76$ \\
Output Noise Power & $\sigma_{out}^2 = \sigma_e^2 \sum_{n} h^2[n]$ \\
\bottomrule
\end{tabular}
\caption{Key quantization formulas.}
\end{table}

\section{Checkpoint Questions}
\textbf{Q1:} An 8-bit ADC has range $\pm 5$V. Calculate step size and quantization noise power. \\
\textbf{A1:} $\Delta = \frac{10}{256} \approx 0.039$ V. $\sigma_e^2 = \frac{(0.039)^2}{12} \approx 1.27 \times 10^{-4}$ W.

\textbf{Q2:} Derive the max amplitude $A$ for a 12-bit ADC (range $\pm 2.5$V) to get SQNR = 70 dB. \\
\textbf{A2:} $\Delta = 5/4096 = 0.00122$ V. $P_e = \frac{\Delta^2}{12} = 1.24 \times 10^{-7}$ W. For 70 dB, $P_x/P_e = 10^7$. $P_x = 1.24$ W. $P_x = A^2/2 \Rightarrow A \approx 1.57$ V.

\textbf{Q3:} A first-order IIR $y[n] = x[n] + 0.8 y[n-1]$ has multiplication noise $e[n]$. Total output noise power? \\
\textbf{A3:} $h_{noise}[n] = (0.8)^n u[n]$. $\sigma_{out}^2 = \sigma_e^2 \sum (0.8^{2n}) = \sigma_e^2 \sum (0.64)^n = \sigma_e^2 \frac{1}{1-0.64} \approx 2.78 \sigma_e^2$.

\end{document}
